\newpage
\section{First exercise}

\textbf{(a)} Initially the code had a problem that it uses the \texttt{NINT} function and not the \texttt{FLOOR} function. This gives incorrect interpolated values when $x_u$ has a decimal part above 0.5. The code was thus changed to 
\begin{verbatim}
    SUBROUTINE interp(xu, T, Tint)
        IMPLICIT NONE
        REAL, INTENT(IN) :: xu
        REAL, DIMENSION(:), INTENT(IN) :: T
        REAL, INTENT(OUT) :: Tint
        REAL :: ax
        INTEGER :: ip, im

        ip = FLOOR(xu) + 1
        im = FLOOR(xu)
        ax = REAL(ip) - xu

        Tint = T(im)*ax + T(ip)*(1 - ax)
\end{verbatim}
Using this corrected code the interpolated temperature at position $x_u=10.4$ was calculated to be $T_{int}=14$. 

\textbf{(b)} However, the code provided had one more issue, since does not use a staggered grid. If we assume a staggered grid then the temperature points a eat the h-points not the u-points($x_u$). are not at the same positions but of set by $\Delta x$. In the staggered grid the we have that $h_j=u_{j-\frac{1}{2}}$. Another way of stating this is that the point $u_j$ is surrounded by 
\begin{equation}
    h_{j}< u_j<u_{j+1}. 
\end{equation}
However, one must note that the staggering can either be forward or backward, however the common way is to use what i showed above. One can thus modify the code so that the position $x_u=10.4$ lies in the position $x_h=9.9$. One thus has the code as:
\begin{verbatim}
    SUBROUTINE interp(xu, T, Tint)
        IMPLICIT NONE
        REAL, INTENT(IN) :: xu
        REAL, DIMENSION(:), INTENT(IN) :: T
        REAL, INTENT(OUT) :: Tint
        REAL :: ax, xh
        INTEGER :: ip, im

        xh = xu-0.5 ! Here we change to h-points

        ip = FLOOR(xh) + 1
        im = FLOOR(xh)
        ax = REAL(ip) - xh

        Tint = T(im)*ax + T(ip)*(1 - ax)
\end{verbatim}

\begin{answer}
    One thus obtains an answer of $T(x_u=10.4)=T(x_h=9.9)=9.5$.
\end{answer}

\textbf{(c)} If an unstaggered grid was used then the original code would work just fine since this would imply that $x_u=x_h$. 